\documentclass[aspectratio=169,11pt]{beamer}

\usepackage{iftex}
\ifXeTeX
  \usepackage{fontspec}
  \usepackage{xeCJK}
  \setmainfont{Times New Roman}
  \setCJKmainfont{SimSun}
\fi
\usepackage{booktabs}
\usepackage{array}
\usepackage{siunitx}
\usepackage{xurl}

\usetheme{Madrid}
\usecolortheme{default}
\setbeamertemplate{navigation symbols}{}
\setbeamertemplate{footline}[frame number]

\newcommand{\movetag}[2]{\textbf{Move #1:} #2}
\newcommand{\steptag}[2]{\textbf{Step #1:} #2}
\newcommand{\claim}[1]{\textcolor{blue!65!black}{#1}}
\newcommand{\gaptext}[1]{\textcolor{red!65!black}{#1}}
\newcommand{\solutiontext}[1]{\textcolor{green!45!black}{#1}}
\newcommand{\src}[1]{\textsuperscript{\scriptsize[\texttt{\detokenize{#1}}]}}
\newcolumntype{L}[1]{>{\raggedright\arraybackslash}p{#1}}

\title[MICE Introduction Map]{Introduction 的 MICE 结构映射}
\subtitle{从当前情境到本文 surrogate workflow}
\author{Linan Jiang}
\institute{Aalto University}
\date{7 April 2026}

\begin{document}

\begin{frame}
  \titlepage
\end{frame}

\begin{frame}{MICE 模型：工程硕士论文引言的四个 Move}
\small
\begin{tabular}{L{0.18\linewidth}L{0.34\linewidth}L{0.38\linewidth}}
\toprule
Move & 功能 & 本文对应问题 \\
\midrule
Move 1 & Current situation & 为什么船用多燃料发动机中的 pilot spray 值得研究？ \\
Move 2 & Problems & 现有试验、CFD、经验模型和图像测量缺少什么？ \\
Move 3 & New approach & 为什么从 Mie 视频学习 surrogate 是合理路线？ \\
Move 4 & Your solution & 本论文具体做什么、怎么做、做到什么范围？ \\
\bottomrule
\end{tabular}

\vspace{0.7em}
\footnotesize
MICE 强调 general-to-specific 和 problem-solution 顺序；本 slide deck 用它检查 introduction 的论证链。
\end{frame}

\begin{frame}{Introduction 的总路线}
\small
\begin{enumerate}
  \item \claim{行业背景}：航运减排和替代燃料推动多燃料船用发动机发展。
  \item \claim{技术背景}：多燃料发动机仍依赖 pilot fuel 实现可靠点火。
  \item \gaptext{工程风险}：pilot diesel spray 的空间发展可能带来撞壁、润滑油稀释和燃烧恶化。
  \item \gaptext{方法缺口}：发动机试验和 CFD 成本高，纯图像测量又不足以形成可快速筛选的预测模型。
  \item \solutiontext{新路线}：用 Mie 高速成像数据构造 reduced-order penetration targets，并训练不确定性感知 surrogate。
\end{enumerate}
\end{frame}

\section{Move 1}

\begin{frame}{Move 1-1：Making a centrality claim}
\small
\movetag{1}{Current situation}

\vspace{0.5em}
\steptag{1}{Making a centrality claim}

\vspace{0.8em}
\begin{block}{本文论证}
国际航运正同时面临温室气体减排、燃料多样化和运行可靠性压力；IMO 2023 战略和行业订单趋势使替代燃料推进系统成为现实工程问题 \src{imo2023ghgstrategy; classnk2025alternativefuels}。
\end{block}

\begin{itemize}
  \item 作用：把研究放在船用能源转型的大背景中。
  \item 关键词：importance, relevance, application.
  \item 引言位置：Research relevance 第 1 段。
\end{itemize}
\end{frame}

\begin{frame}{Move 1-2：Making topic generalizations}
\small
\movetag{1}{Current situation}

\vspace{0.5em}
\steptag{2}{Making topic generalizations}

\vspace{0.8em}
\begin{block}{本文论证}
双燃料和多燃料船用发动机不仅需要替代燃料供应，也需要可靠点火；低压燃气、甲醇和氨方案通常仍需要少量 pilot fuel \src{wingd_xdf_manual_2021; wingd_lng_faq_2025; classnk_methanol_2024; wingd_ammonia_faq_2025}。
\end{block}

\begin{itemize}
  \item 作用：解释 pilot fuel 在系统中的一般功能。
  \item 进一步收窄：pilot diesel spray 不只是 ignition aid，其空间发展也影响工程风险。
  \item 引言位置：Research relevance 第 2--3 段。
\end{itemize}
\end{frame}

\begin{frame}{Move 1-3：Reviewing previous research and solutions}
\small
\movetag{1}{Current situation}

\vspace{0.5em}
\steptag{3}{Reviewing previous research and solutions}

\vspace{0.8em}
\begin{block}{本文论证}
喷雾与撞壁风险通常通过发动机试验、CFD/喷雾燃烧仿真和光学喷雾实验评估；柴油喷雾研究已经形成了高速成像、经验关联式、现象学模型和 CFD 等方法基础 \src{Linne2013SprayImaging; HiroyasuArai1990; NaberSiebers1996; liu2020sprayreview}。
\end{block}

\begin{itemize}
  \item 作用：承认领域已有成熟基础，避免把问题写成“没人研究过喷雾”。
  \item 引言位置：Current engineering practice 和 Existing solutions。
\end{itemize}
\end{frame}

\section{Move 2}

\begin{frame}{Move 2-1：Indicating weaknesses}
\small
\movetag{2}{Problems}

\vspace{0.5em}
\steptag{1}{Indicating weaknesses}

\vspace{0.8em}
\begin{itemize}
  \item 发动机试验接近真实工况，但昂贵、慢，难以隔离单一设计因素。
  \item 高保真 CFD 信息丰富，但需要复杂模型、大量输入假设和显著计算资源。
  \item 一个 reacting diesel-spray LES 例子需要约 48,000 CPU core-hours，并且可达 20 天 wall-clock，只覆盖 SOI 后 \(\SI{2}{ms}\) \src{CurtisNiemeyerSung2017}。
  \item 经验和现象学模型紧凑可解释，但对短脉冲、强瞬态和多次喷射情形不够稳健。
  \item 光学边界具有方法依赖性，不等于隐藏物理真值 \src{ecn_liquid_penetration; ecn_liquid_penetration_accessory}。
\end{itemize}
\end{frame}

\begin{frame}{Move 2-2：Identifying a gap}
\small
\movetag{2}{Problems}

\vspace{0.5em}
\steptag{2}{Identifying a gap}

\vspace{0.8em}
\begin{block}{本文 gap statement}
缺少一种快速、可重复的方法，用于预测预喷柴油非反应液相包络随时间的演化，从而为后续撞壁风险分析提供有依据的几何输入。
\end{block}

\begin{itemize}
  \item 不是缺少喷雾研究本身。
  \item 缺的是光学实验和高成本验证之间的 intermediate-layer predictor。
  \item 目标不是替代 CFD 或发动机试验，而是在它们之前缩小候选空间。
\end{itemize}
\end{frame}

\section{Move 3}

\begin{frame}{Move 3-1：Introducing a potential solution}
\small
\movetag{3}{Proposing a new approach}

\vspace{0.5em}
\steptag{1}{Introducing a potential solution}

\vspace{0.8em}
\begin{block}{本文新路线}
直接从大规模 Mie-scattering optical spray data 中学习一个 condition-parametric surrogate，把测得的 penetration trajectories 转换成可快速查询的预测模型。
\end{block}

\begin{itemize}
  \item Mie 视频已经捕获液相包络行为。
  \item Reduced-order fitting 把原始观测转换为稳定学习目标。
  \item Surrogate 模型位于 raw measurement 和 physics-based simulation 之间。
\end{itemize}
\end{frame}

\begin{frame}{Move 3-2：Justifying the approach}
\small
\movetag{3}{Proposing a new approach}

\vspace{0.5em}
\steptag{2}{Justifying the approach}

\vspace{0.8em}
\begin{itemize}
  \item \solutiontext{保留实验保真度}：训练目标来自实际 Mie 成像和穿深轨迹。
  \item \solutiontext{提高筛选速度}：避免每个候选工况都进入 CFD 或发动机试验。
  \item \solutiontext{适应数据扩展}：条件到轨迹参数的映射可随数据集更新。
  \item \solutiontext{可解释和可审计}：标定、边界定义、轨迹拟合和不确定性回归都有中间产物。
  \item \solutiontext{符合科学图像处理实践}：Fiji/ImageJ、scikit-image、OpenPIV、PIVlab 等生态强调脚本化、可检查、可复用流程 \src{Schindelin2012Fiji; VanDerWalt2014ScikitImage; OpenPIV; ThielickeSonntag2021PIVlab}。
\end{itemize}
\end{frame}

\section{Move 4}

\begin{frame}{Move 4-1：Research aims}
\small
\movetag{4}{Your solution}

\vspace{0.5em}
\steptag{1}{Research aims}

\vspace{0.8em}
\begin{block}{本文目标}
开发一个用于预喷柴油的、数据驱动的非反应液相包络预测模型，从喷射工况与时间信息预测喷雾包络和穿深趋势，并给出适于工程解释的不确定性描述。
\end{block}

\begin{itemize}
  \item 输出：mean penetration trend + uncertainty.
  \item 用途：船用双燃料发动机开发中的早期撞壁风险筛查。
\end{itemize}
\end{frame}

\begin{frame}{Move 4-2：Methods}
\small
\movetag{4}{Your solution}

\vspace{0.5em}
\steptag{2}{Methods}

\vspace{0.8em}
\begin{enumerate}
  \item 室温加压空气舱中的柴油喷雾高速 Mie 成像。
  \item 人工标定喷嘴中心和像素尺度。
  \item CUDA 加速批处理提取 plume-resolved temporal observables 和边界几何量。
  \item 降阶拟合构造 penetration targets，并处理右删失与时间对齐。
  \item 多阶段学习：Stage-1 MSE、Stage-2 heteroscedastic NLL、teacher-guided refinement。
  \item 将预测均值和不确定性接入简化几何风险筛查。
\end{enumerate}
\end{frame}

\begin{frame}{Move 4-3：Thesis scope}
\small
\movetag{4}{Your solution}

\vspace{0.5em}
\steptag{3}{Thesis scope}

\vspace{0.8em}
\begin{itemize}
  \item 范围内：受控加压舱、室温、非反应、非蒸发 pilot diesel spray。
  \item 范围内：光学布置下观察到的宏观液相包络几何及其时间演化。
  \item 范围外：燃烧化学、蒸发、双燃料点火耦合、真实热壁面相互作用、真实发动机 wall-film truth。
  \item 模型边界：主要面向数据覆盖范围内的插值预测；明显 OOD 外推不作同等级保证。
\end{itemize}
\end{frame}

\begin{frame}{Move 4-4：Main outcome}
\small
\movetag{4}{Your solution}

\vspace{0.5em}
\steptag{4}{Describing the main outcome}

\vspace{0.8em}
\begin{enumerate}
  \item TB 级 Mie 视频到 plume-resolved observables 的 CUDA 加速工作流。
  \item 面向轨迹对齐、稳健外推和右删失缓解的 reduced-order fitting。
  \item 结合 MSE、heteroscedastic NLL 和 teacher-guided refinement 的多阶段 surrogate。
  \item 将 penetration prediction 与 downstream geometric collision evaluation 解耦的概率化筛查原型。
\end{enumerate}
\end{frame}

\begin{frame}{Move 4-5：Thesis structure}
\small
\movetag{4}{Your solution}

\vspace{0.5em}
\steptag{5}{Overview of thesis structure}

\vspace{0.8em}
\begin{tabular}{L{0.25\linewidth}L{0.62\linewidth}}
\toprule
部分 & 作用 \\
\midrule
文献综述与技术依据 & 放置于柴油喷雾成像、光学诊断和 surrogate 建模背景下。 \\
材料与方法 & 实验数据、图像处理、特征提取、轨迹拟合和模型训练。 \\
结果 & 报告已归档的处理资产和模型证据。 \\
结论 & 总结结果、限制和后续开发需求。 \\
\bottomrule
\end{tabular}
\end{frame}

\begin{frame}{Move/Step 检查表}
\small
\begin{tabular}{L{0.19\linewidth}L{0.30\linewidth}L{0.40\linewidth}}
\toprule
MICE 位置 & Introduction 小节 & 核心论证 \\
\midrule
Move 1-1 & Research relevance & 航运减排和替代燃料使问题重要。 \\
Move 1-2 & Research relevance & 多燃料发动机依赖 pilot fuel，喷雾空间发展有风险。 \\
Move 1-3 & Current practice / Existing solutions & 试验、CFD、光学诊断和经验模型已有基础。 \\
Move 2-1 & Existing solutions & 各类方法在成本、瞬态、可解释性或真值定义上有弱点。 \\
Move 2-2 & Problem statement & 缺少光学实验和高成本验证之间的快速中间层预测器。 \\
Move 3-1 & Problem statement 后段 & 从 Mie 数据学习 surrogate。 \\
Move 3-2 & Why workflow necessary & 降阶目标、审计性、不确定性和高通量处理使路线合理。 \\
Move 4-1--5 & Aim / Methods / Scope / Contributions / Structure & 本论文具体方案、边界、结果和章节组织。 \\
\bottomrule
\end{tabular}
\end{frame}

\begin{frame}{Source Keys}
\small
本 slide deck 使用论文 bibliography 中的 citation keys 作轻量标注；正式参考文献仍以 thesis PDF 的 bibliography 为准。

\vspace{0.7em}
\begin{tabular}{L{0.36\linewidth}L{0.52\linewidth}}
\toprule
主题 & source keys \\
\midrule
航运与替代燃料 & imo2023ghgstrategy; classnk2025alternativefuels \\
pilot fuel 背景 & wingd\_xdf\_manual\_2021; wingd\_lng\_faq\_2025; classnk\_methanol\_2024; wingd\_ammonia\_faq\_2025 \\
喷雾实验与模型 & Linne2013SprayImaging; HiroyasuArai1990; NaberSiebers1996; liu2020sprayreview \\
CFD 成本 & CurtisNiemeyerSung2017 \\
光学边界定义 & ecn\_liquid\_penetration; ecn\_liquid\_penetration\_accessory \\
图像处理生态 & Schindelin2012Fiji; VanDerWalt2014ScikitImage; OpenPIV; ThielickeSonntag2021PIVlab \\
\bottomrule
\end{tabular}
\end{frame}

\end{document}
